\SourceLeaf{361}%
calculation is to be carried out. Strictly speaking, all the other ways of solving the problem, including \textsc{Bessel}'s in particular, also have this character of indirect operations in part. What matters most, however, is that only the last of the successive approximations has to be kept. It contains the entire substance of the calculation, with a brevity and clarity that no other method offers.

The formulas for solving the problem on the spheroid differ from those for the sphere only in that certain intermediate quantities, which are constant in the latter. depend on latitude in the former. These quantities can therefore be entered in an auxiliary table with latitude as its argument. With such a table at hand, the calculation on the spheroid is just as easy in every specific case as it is on the sphere. Such a table is appended to the memoir for the zone from 51---54 degrees of latitude, which suffices for the Hanoverian triangulation network. It uses the value of the flattening that \textsc{Bessel} derived from all the degree measurements made up to that time, which also served as the basis of the first memoir. Anyone wishing to apply the same method to a triangulation network outside this zone would first have to extend the table as far as needed, or construct a new table if a different flattening were to be used. In working up a large triangulation network, this preliminary work is of no consequence, and the convenience in use amply repays the effort spent on it. Where only one or a couple of specific solutions to the problem in question are wanted, however, the method was not intended particularly for that case.
